\section{Discussion}
\label{sec:discussion}

\subsection{Key Findings}

\textbf{Finding 1:} The EDMP pipeline is computationally tractable and produces statistically significant domain differences (ANOVA F=1242.59, p$<$0.001) with perfect reproducibility. The pipeline completes in under 30 minutes on a single GPU.

\textbf{Finding 2:} Synthetic data produces insufficient variance for domain discrimination. This is a data quality issue, not a fundamental method flaw.

\subsection{Limitations}

\textbf{L1: Only existence validated, not predictive power.} We tested whether EDMP scores CAN be computed, not whether they PREDICT training utility. The core hypothesis remains untested.

\textbf{L2: Synthetic data used instead of real Pile domains.} Due to data access constraints, we used synthetic domain-distinguishable text lacking vocabulary diversity.

\textbf{L3: Single embedder evaluated.} Rankings might differ with BGE-large, OpenAI embeddings, or other E5 variants.

\textbf{L4: Main predictions (P1, P2, P3) untested.} None of the three main predictions were tested; these depend on completing the causal mechanism chain.

\subsection{Broader Impact}

If validated, EDMP could democratize LLM data optimization by eliminating expensive proxy model training. Training-free optimization methods also reduce energy consumption compared to iterative training approaches.
